% Code settings
\lstset{
  language=Python,
  basicstyle=\ttfamily\small,
  numbers=left,
  numberstyle=\footnotesize,
  stepnumber=1,
  numbersep=2.0mm,
  breaklines=true}

\chapter{Codes}\label{ch:codes}

\section{Device Metrics JSONL Consolidation}
Listing~\ref{lst:jsonl} shows the core merge logic from \texttt{Shared\_pipeline\_Files/tools/jsonl\_to\_json.py}, which folds the append-only on-device log into the legacy \texttt{all\_scan\_metrics.json} format consumed by plotting scripts.

\lstinputlisting[caption={JSONL-to-JSON metrics consolidation (excerpt)},label={lst:jsonl}]{codes/jsonl_to_json_snippet.py}

\section{Composite Feasibility Score}
The plotting pipeline ranks models using a weighted composite of normalized quality, cascade tier priority, and hardware cost:
\begin{lstlisting}[caption={Feasibility composite (from device\_metrics\_lib.py)},label={lst:composite}]
composite = 0.45 * quality_norm + 0.30 * tier_weight + 0.25 * (1 - cost_norm)
cost_score = stage_total_ms + 10 * mem_mb
quality = mean(accuracy, f1, roc_auc)
\end{lstlisting}

\endinput
